\documentclass[10pt,conference]{IEEEtran}
\usepackage{amsmath}
\usepackage{booktabs}

\title{HELIOS: A Multi-Language Framework for Unified Signal Processing, Information Theory, and Stochastic Systems}
\author{George David Tsitlauri\\
\textit{Dept. of Informatics \& Telecommunications}\\
\textit{University of Thessaly, Greece}\\
gdtsitlauri@gmail.com}
\date{2026}

\begin{document}
\maketitle

\begin{abstract}
HELIOS is an open-source framework that unifies classical digital signal processing, information theory, stochastic systems, and neural temporal prediction under a single research workflow. The proposed HELIOS-SPECTRUM pipeline combines wavelet decomposition, stochastic regime characterization, information-theoretic feature selection, and temporal convolutional prediction in a single causal analysis stack. The strongest empirical evidence in the current release is concentrated in the
HELIOS-SPECTRUM benchmark suite, where the proposed pipeline outperforms the
committed FFT+ML, ARIMA, and CNN baselines across ECG, seismic, financial, and
audio datasets, while the surrounding DSP, coding, and stochastic modules serve
as supporting components of the broader signal-research platform.
\end{abstract}

\section{Introduction}
HELIOS targets reproducible signal research across Python, Julia, R, and Octave-compatible workflows.

\section{DSP Foundation}
The framework includes FFT, STFT, filter design, and wavelet denoising APIs with Julia bridge support.

\section{Information Theory}
Entropy, mutual information, source coding, arithmetic coding, Gaussian rate-distortion baselines, and an expanded channel coding stack including Hamming, convolutional/Viterbi, LDPC-style, and turbo-style baselines are implemented.

\section{Stochastic Processes}
The stochastic module supports Markov estimation, hitting-time and absorption analysis, Baum-Welch-style HMM updates, AR-style forecasting, Viterbi decoding, continuous-process simulation, and Kalman tracking.

\section{HELIOS-SPECTRUM Algorithm}
HELIOS-SPECTRUM decomposes signals into sub-bands, scores them via mutual information, predicts targets with a compact temporal convolutional model, and estimates directed dependencies via Granger-style lag scoring.

\section{Multi-Language Architecture}
Python acts as the orchestrator while Julia, R, and Octave bridges activate when the runtimes are installed.

\section{Experiments}
% FFT Benchmark Table (inlined from results/dsp/fft_benchmarks.csv)
\begin{table}[h]
\centering
\caption{FFT Benchmark}
\begin{tabular}{lcc}
\toprule
Backend & Seconds & Speedup \\
\midrule
NumPy & 0.00030 & 1.00 \\
Julia & 0.00011 & 2.25 \\
\bottomrule
\end{tabular}
\end{table}
% Channel Coding Table (inlined from results/information_theory/channel_coding_ber.csv)
\begin{table}[h]
\centering
\caption{Channel Coding at 4 dB}
\begin{tabular}{lc}
\toprule
Code & BER \\
\midrule
Hamming (7,4) & 0.1128 \\
Hamming (15,11) & 0.0549 \\
Convolutional + Viterbi & 0.00015 \\
Reed-Solomon-like & 0.0152 \\
LDPC (bitflip) & 0.00006 \\
Turbo (simplified) & 0.000015 \\
Neural Autoencoder & 0.000037 \\
\bottomrule
\end{tabular}
\end{table}
% Compression Ratio Table (inlined from results/information_theory/compression_comparison.csv)
\begin{table}[h]
\centering
\caption{Compression Ratio}
\begin{tabular}{lcc}
\toprule
Algorithm & Ratio & Shannon Limit \\
\midrule
Huffman & 2.07 & 2.00 \\
LZ78 & 1.00 & 2.00 \\
Arithmetic & 13.76 & 2.00 \\
\bottomrule
\end{tabular}
\end{table}
% HELIOS-SPECTRUM vs Baselines Table (inlined from results/helios_spectrum/comparison_table.csv)
\begin{table}[h]
\centering
\caption{HELIOS-SPECTRUM vs Baselines}
\begin{tabular}{lc}
\toprule
Dataset & HELIOS-SPECTRUM & FFT+ML & ARIMA & CNN \\
\midrule
ECG & 0.0411 & 0.0574 & 0.062 & 0.047 \\
Seismic & 0.331 & 0.905 & 0.978 & 0.381 \\
Financial & 0.651 & 4.53 & 4.90 & 0.749 \\
Audio & 0.248 & 0.787 & 0.850 & 0.285 \\
\bottomrule
\end{tabular}
\end{table}

\subsection{Interpretation Boundary}
The current HELIOS release should be read as a unified signal-research platform
with one especially clear empirical center: HELIOS-SPECTRUM. The FFT, channel
coding, compression, and stochastic modules provide meaningful supporting
artifacts and demonstrate breadth, but they are not all intended to carry equal
weight as standalone benchmark papers. The most defensible flagship claim in
the present repository is that HELIOS-SPECTRUM provides a reproducible,
multi-language signal-analysis pipeline with strong comparative results on the
committed benchmark families.


\section{Conclusion}
The current release establishes a reproducible, extensible foundation for the
full HELIOS research agenda while already providing a strong empirical core in
HELIOS-SPECTRUM and a credible set of supporting DSP, coding, and stochastic
artifacts around that center.

\end{document}
